%% ----------------------------------------------------------------------------
%%  §1.2  Alcune conseguenze degli assiomi dei numeri reali
%% ----------------------------------------------------------------------------
\section{Alcune conseguenze degli assiomi}
\label{sec:conseguenze}

Nel paragrafo precedente sono state elencate le proprietà dei numeri reali che vengono
assunte come assiomi. \emph{Tutte} le altre proprietà e tutti i teoremi esposti in
questo libro discendono dagli assiomi. Sono conseguenze degli assiomi anche quelle
proprietà elementari che fanno parte del «bagaglio matematico» di ogni studente, come il
fatto che \emph{un prodotto è nullo quando almeno uno dei due fattori è nullo}, oppure la
regola dei segni per il prodotto (che dagli studenti delle scuole elementari è talvolta
accettata come imposizione, perché incompresa) che schematicamente si enuncia:
\emph{meno per meno fa più}; oppure la norma, di frequente applicazione nel risolvere le
disequazioni, secondo cui \emph{moltiplicando entrambi i membri per una quantità
negativa il verso della disequazione cambia}.

Di seguito dimostriamo alcune di queste proprietà. Il lettore non si lasci ingannare
dalla loro apparente banalità: il valore di questo paragrafo sta nel \emph{metodo}. Ogni
passaggio deve essere giustificato da un assioma o da una proprietà già dimostrata, e
nient'altro; è lo stesso rigore che useremo per i teoremi più difficili del libro, e
conviene allenarlo su enunciati di cui si conosce già la risposta.

\subsection{Conseguenze degli assiomi delle operazioni}

\begin{proposizione}[Legge di semplificazione per la somma]\label{prop:semplificazione-somma}
Siano $a,b,c\in\R$. Se $a+b=a+c$, allora $b=c$.
\end{proposizione}

\begin{schema}
\Ipotesi{$a,b,c\in\R$ e $a+b=a+c$.}
\Tesi{$b=c$.}
\end{schema}

\begin{strategia}
Vogliamo «far sparire» $a$ dall'uguaglianza $a+b=a+c$. L'unico strumento che gli assiomi
ci offrono per cancellare un addendo è l'opposto: per \ref{ass:opposto} esiste $-a$ con
$a+(-a)=0$. L'idea è dunque di partire da $b$, scriverlo come $0+b$, sostituire $0$ con
$(-a)+a$ e poi, usando la proprietà associativa, far comparire il blocco $a+b$, che per
ipotesi vale $a+c$. A quel punto si ripercorre la strada a ritroso e si arriva a $c$.
\end{strategia}

\begin{dimostrazione}
Utilizziamo gli assiomi \ref{ass:neutri} e \ref{ass:opposto} (esistenza dello zero e
dell'opposto $-a$) e le proprietà commutativa \ref{ass:commutativa} e associativa
\ref{ass:associativa}:
\[
  b = 0+b = \bigl[a+(-a)\bigr]+b = \bigl[(-a)+a\bigr]+b = (-a)+(a+b).
\]
Essendo per ipotesi $a+b=a+c$, si ottiene
\begin{align*}
  b &= (-a)+(a+b) = (-a)+(a+c) = \bigl[(-a)+a\bigr]+c \\
    &= \bigl[a+(-a)\bigr]+c = 0+c = c+0 = c,
\end{align*}
che è la tesi.
\end{dimostrazione}

\begin{proposizione}[Legge di semplificazione per il prodotto]\label{prop:semplificazione-prodotto}
Siano $a,b,c\in\R$. Se $a\cdot b=a\cdot c$ e $a\neq0$, allora $b=c$.
\end{proposizione}

\begin{schema}
\Ipotesi{$a,b,c\in\R$, \ $a\cdot b=a\cdot c$ \ e \ $a\neq0$.}
\Tesi{$b=c$.}
\end{schema}

\begin{strategia}
Si procede esattamente come nella dimostrazione precedente, scambiando la somma con il
prodotto, lo zero con l'uno e l'opposto con l'inverso. L'unica accortezza è ricordare
che l'inverso $a^{-1}$ esiste \emph{purché} $a\neq0$, per \ref{ass:inverso}: ecco
perché l'ipotesi $a\neq0$ è indispensabile.
\end{strategia}

\begin{dimostrazione}
Poiché $a\neq0$, esiste l'inverso $a^{-1}$. Allora
\begin{align*}
  b &= b\cdot1 = 1\cdot b = (a\cdot a^{-1})\cdot b = (a^{-1}\cdot a)\cdot b \\
    &= a^{-1}\cdot(a\cdot b) = a^{-1}\cdot(a\cdot c) = (a^{-1}\cdot a)\cdot c \\
    &= (a\cdot a^{-1})\cdot c = 1\cdot c = c\cdot 1 = c. \qedhere
\end{align*}
\end{dimostrazione}

\begin{proposizione}[Legge di annullamento del prodotto]\label{prop:annullamento}
\index{legge di annullamento del prodotto}
Siano $a,b\in\R$. Il prodotto $a\cdot b$ è nullo se e soltanto se almeno uno dei due
fattori è nullo.
\end{proposizione}

\begin{schema}
\Ipotesi{$a,b\in\R$.}
\Tesi{$a\cdot b=0 \ \sse\ (a=0 \text{ oppure } b=0)$.}
\end{schema}

\begin{strategia}
L'enunciato è una doppia implicazione («se e soltanto se»): dobbiamo dimostrare due
cose distinte.
\begin{itemize}
  \item \textbf{Il «se»:} se $a=0$ oppure $b=0$ allora $a\cdot b=0$. Basta provare che
  $a\cdot0=0$ per ogni $a$. Sembra ovvio, ma non è un assioma! Lo zero è definito come
  elemento neutro della \emph{somma}, e nulla dice su come si comporta nel
  \emph{prodotto}. Il trucco è scrivere $0=1+0$ e usare la proprietà distributiva.
  \item \textbf{Il «soltanto se»:} se $a\cdot b=0$ allora $a=0$ oppure $b=0$. Se $a=0$
  non c'è nulla da dimostrare; se $a\neq0$ possiamo moltiplicare per $a^{-1}$ e usare la
  prima parte.
\end{itemize}
\end{strategia}

\begin{dimostrazione}
Proviamo preliminarmente l'implicazione con il «se», cioè che $a\cdot0=0$ per ogni
numero reale $a$. Ricordiamo che lo zero è, per l'assioma \ref{ass:neutri}, l'elemento
neutro rispetto alla somma, cioè $a+0=a$ per ogni $a$; e che $1$, elemento neutro
rispetto al prodotto, soddisfa $a\cdot1=a$. In base alla proprietà distributiva
\ref{ass:distributiva} abbiamo allora
\[
  a + a\cdot0 = a\cdot1 + a\cdot0 = a\cdot(1+0) = a\cdot1 = a = a+0,
\]
da cui, semplificando entrambi i membri in base alla
\cref{prop:semplificazione-somma}, otteniamo $a\cdot0=0$. Per la proprietà commutativa
vale anche $0\cdot b=0$.

Proviamo ora l'implicazione con il «soltanto se»: supponiamo $a\cdot b=0$. Se $a=0$
la tesi è raggiunta. Altrimenti, se $a\neq0$, esiste l'inverso $a^{-1}$ e si ha
\[
  b = b\cdot1 = b\cdot(a\cdot a^{-1}) = (b\cdot a)\cdot a^{-1}
    = (a\cdot b)\cdot a^{-1} = 0\cdot a^{-1} = 0 .
\]
Si noti che nell'ultimo passaggio abbiamo utilizzato quanto già provato nella prima
parte della dimostrazione.
\end{dimostrazione}

\begin{osservazione}[Perché non si può dividere per zero]
La \cref{prop:annullamento} spiega perché nell'ambito dei numeri reali non sia
possibile la divisione per zero, cioè perché nell'assioma \ref{ass:inverso} si richieda
$a\neq0$. Infatti, se fosse $a=0$, allora $a\cdot b=0\cdot b=0$ per ogni numero reale
$b$, e perciò \emph{non esiste} alcun numero reale $0^{-1}$ tale che $0\cdot0^{-1}=1$:
se esistesse, avremmo $1=0\cdot0^{-1}=0$, contro l'assioma \ref{ass:neutri} che
richiede $0\neq1$. La divisione per zero non è «vietata» da una regola arbitraria: è
\emph{impossibile} perché contraddittoria con gli assiomi.
\end{osservazione}

\begin{proposizione}[Unicità dell'opposto e dell'inverso]\label{prop:unicita}
\leavevmode
\begin{enumerate}
  \item L'opposto di un numero reale è unico.
  \item L'inverso di un numero reale non nullo è unico.
  \item Per ogni $a\in\R$ vale $-(-a)=a$.
\end{enumerate}
\end{proposizione}

\begin{schema}
\Ipotesi{$a\in\R$; per (ii) $a\neq0$.}
\Tesi{(i) se $a+x=0$ e $a+y=0$, allora $x=y$; \ (ii) se $a\cdot x=1$ e $a\cdot y=1$,
allora $x=y$; \ (iii) $-(-a)=a$.}
\end{schema}

\begin{strategia}
Per dimostrare che un oggetto è unico si procede sempre allo stesso modo: si suppone
che ce ne siano due, $x$ e $y$, con la stessa proprietà, e si dimostra che $x=y$. Qui
la proprietà è «$a+x=0$», e la legge di semplificazione fa il resto. Il punto (iii) è
un'applicazione dell'unicità: $-(-a)$ è \emph{per definizione} l'opposto di $-a$; ma
anche $a$ è un opposto di $-a$, perché $(-a)+a=0$; dunque i due coincidono.
\end{strategia}

\begin{dimostrazione}
(i) In base all'assioma \ref{ass:opposto}, per ogni numero reale $a$ esiste l'opposto
$-a$, tale che $a+(-a)=0$. Se supponiamo che risulti anche $a+b=0$ per qualche $b$,
allora $a+(-a)=a+b$ e, per la legge di semplificazione (\cref{prop:semplificazione-somma}),
$-a=b$: l'opposto è unico.

(ii) Stessa dimostrazione, con il prodotto al posto della somma e la
\cref{prop:semplificazione-prodotto} al posto della \cref{prop:semplificazione-somma}.

(iii) Il numero $-(-a)$ è, per definizione, l'opposto di $-a$. Ma essendo
$a+(-a)=(-a)+a=0$, risulta che anche $a$ è un opposto di $-a$; per l'unicità
dell'opposto provata in (i), $a=-(-a)$.
\end{dimostrazione}

\begin{proposizione}[Regola dei segni]\label{prop:segni}\index{regola dei segni}
Per ogni coppia di numeri reali $a$, $b$ risulta
\begin{enumerate}
  \item $(-a)\cdot b = -(a\cdot b)$;
  \item $(-a)\cdot(-b) = a\cdot b$.
\end{enumerate}
\end{proposizione}

\begin{schema}
\Ipotesi{$a,b\in\R$.}
\Tesi{(i) $(-a)\cdot b=-(a\cdot b)$; \quad (ii) $(-a)\cdot(-b)=a\cdot b$.}
\end{schema}

\begin{strategia}
Che cosa significa dimostrare che $(-a)\cdot b$ è uguale a $-(a\cdot b)$? Significa
dimostrare che $(-a)\cdot b$ è \emph{l'opposto} di $a\cdot b$, cioè che sommato ad
$a\cdot b$ dà zero; per l'unicità dell'opposto (\cref{prop:unicita}) ciò basta. E la
somma $(-a)\cdot b+a\cdot b$ si calcola con la proprietà distributiva, raccogliendo
$b$. Il punto (ii), il celebre «meno per meno fa più», si ottiene applicando due
volte il punto (i) e poi la regola $-(-x)=x$.
\end{strategia}

\begin{dimostrazione}
(i) Per la proprietà distributiva \ref{ass:distributiva} e per la
\cref{prop:annullamento} si ha
\[
  (-a)\cdot b + a\cdot b = \bigl[(-a)+a\bigr]\cdot b = 0\cdot b = 0,
\]
da cui $(-a)\cdot b$ è l'opposto di $a\cdot b$, cioè $(-a)\cdot b=-(a\cdot b)$.

(ii) Come conseguenza del punto (i), applicato due volte, e della proprietà commutativa
\ref{ass:commutativa}, abbiamo
\[
  (-a)\cdot(-b) = -\bigl[a\cdot(-b)\bigr] = -\bigl[(-b)\cdot a\bigr]
  = -\bigl[-(b\cdot a)\bigr] = -\bigl[-(a\cdot b)\bigr];
\]
la conclusione segue infine dalla \cref{prop:unicita}(iii), essendo
$-[-(a\cdot b)]=a\cdot b$.
\end{dimostrazione}

\begin{notastorica}[«Meno per meno fa più»]
La regola dei segni fu a lungo un ostacolo concettuale. Nel Cinquecento i numeri
negativi erano chiamati \emph{numeri absurdi} o \emph{ficti} (Cardano, Stifel) e ancora
nel 1770 Eulero, nella sua \emph{Vollständige Anleitung zur Algebra}, giustificava la
regola con un argomento poco convincente: il prodotto $(-a)\cdot(-b)$ può valere
soltanto $+ab$ oppure $-ab$; poiché $-ab$ è già il valore di $(-a)\cdot b$, esso «non
può che essere» $+ab$. Lo scrittore
Stendhal racconta nella \emph{Vita di Henry Brulard} di aver sofferto da ragazzo per
questa regola, che nessuno dei suoi insegnanti sapeva spiegare. La dimostrazione che
abbiamo appena dato mostra che la regola non è una convenzione: è l'\emph{unico} modo
di definire il prodotto di due numeri negativi se si vuole che continui a valere la
proprietà distributiva.
\end{notastorica}

\subsection{Conseguenze degli assiomi dell'ordinamento}

Gli assiomi del \cref{sec:assiomi} relativi all'ordinamento si riferiscono alla
relazione di minore o uguale ($\le$). Ricordiamo che le relazioni $\ge$, $<$, $>$ sono
state definite a partire da essa: $a\ge b$ significa $b\le a$, e $a<b$ significa
«$a\le b$ e $a\neq b$». Pertanto la relazione $\ge$ gode di proprietà analoghe a quelle
di $\le$. Nelle dimostrazioni che seguono useremo in modo sistematico la seguente
riformulazione dell'ordine in termini di segno di una differenza: essa ci permetterà
di ricondurre ogni disuguaglianza a una proprietà dei numeri $\ge0$.

\begin{proposizione}\label{prop:ordine-differenza}
Per ogni $a,b\in\R$ la relazione $a\le b$ è equivalente alla relazione $b-a\ge0$.
\end{proposizione}

\begin{schema}
\Ipotesi{$a,b\in\R$.}
\Tesi{$a\le b \ \sse\ 0\le b-a$.}
\end{schema}

\begin{strategia}
L'unico assioma che lega l'ordine alla somma è \ref{ass:ordine-somma}: «se $a\le b$
allora $a+c\le b+c$». Per passare da $a\le b$ a $0\le b-a$ basta scegliere $c=-a$; per
tornare indietro, basta scegliere $c=a$.
\end{strategia}

\begin{dimostrazione}
Se $a\le b$, allora per \ref{ass:ordine-somma} con $c=-a$:
\[
  a+(-a)\le b+(-a), \qquad\text{cioè}\qquad 0\le b-a .
\]
Viceversa, se $0\le b-a$, sempre per \ref{ass:ordine-somma} con $c=a$, e per la
proprietà associativa dell'addizione,
\[
  a = 0+a \le (b-a)+a = b+\bigl[(-a)+a\bigr] = b+0 = b. \qedhere
\]
\end{dimostrazione}

\begin{proposizione}[Proprietà transitiva]\label{prop:transitiva}
Se $a\le b$ e $b\le c$, allora $a\le c$.
\end{proposizione}

\begin{schema}
\Ipotesi{$a,b,c\in\R$, \ $a\le b$ \ e \ $b\le c$.}
\Tesi{$a\le c$.}
\end{schema}

\begin{strategia}
Traduciamo tutto in termini di numeri non negativi con la
\cref{prop:ordine-differenza}: le ipotesi dicono che $b-a\ge0$ e $c-b\ge0$. L'assioma
\ref{ass:ordine-prodotto} ci dice che la somma di due numeri non negativi è non
negativa; e la somma di $b-a$ e $c-b$ è proprio $c-a$.
\end{strategia}

\begin{dimostrazione}
Supponiamo $a\le b$ e $b\le c$; per la \cref{prop:ordine-differenza} risulta
$0\le b-a$ e $0\le c-b$. Dall'assioma \ref{ass:ordine-prodotto} si ottiene
\[
  0 \le (b-a)+(c-b) = c-a,
\]
che, ancora per la \cref{prop:ordine-differenza}, equivale ad $a\le c$.
\end{dimostrazione}

\begin{proposizione}\label{prop:segno-opposto}
Per ogni $a\in\R$ risulta $a\ge0$ se e soltanto se $-a\le0$, e $a\le0$ se e soltanto
se $-a\ge0$. Le stesse equivalenze valgono con le disuguaglianze strette.
\end{proposizione}

\begin{schema}
\Ipotesi{$a\in\R$.}
\Tesi{$a\ge0\sse -a\le0$; \quad $a\le0\sse -a\ge0$; \quad analogamente con $>$ e
$<$.}
\end{schema}

\begin{strategia}
Sommare uno stesso numero ai due membri conserva la disuguaglianza
(\ref{ass:ordine-somma}): per passare da $0\le a$ a $-a\le0$ si somma $-a$, per tornare
indietro si somma $a$. Per le disuguaglianze strette basta osservare in più che $a=0$ se
e soltanto se $-a=0$.
\end{strategia}

\begin{dimostrazione}
Per \ref{ass:ordine-somma}, se $0\le a$ allora $0+(-a)\le a+(-a)$, cioè $-a\le0$.
Viceversa, se $-a\le0$ allora, sommando $a$ ad ambo i membri, $(-a)+a\le0+a$, cioè
$0\le a$. Allo stesso modo, se $a\le0$, sommando $-a$ si ottiene $0\le-a$; e se
$0\le-a$, sommando $a$ si ottiene $a\le0$. Infine $-0=0$, perché $0+0=0$ e l'opposto è
unico (\cref{prop:unicita}); e se $-a=0$, allora $a=-(-a)=-0=0$. Dunque $a=0$ se e
soltanto se $-a=0$, e le equivalenze valgono anche con le disuguaglianze strette.
\end{dimostrazione}

\begin{proposizione}[Prodotto di una disuguaglianza per un numero]\label{prop:ordine-prodotto}
\leavevmode
\begin{enumerate}
  \item Se $a\le b$ e $c\ge0$, allora $a\cdot c\le b\cdot c$.
  \item Se $a\le b$ e $c\le0$, allora $a\cdot c\ge b\cdot c$.
\end{enumerate}
\end{proposizione}

\begin{schema}
\Ipotesi{$a,b,c\in\R$ e $a\le b$; nel caso (i) $c\ge0$, nel caso (ii) $c\le0$.}
\Tesi{(i) $ac\le bc$; \quad (ii) $ac\ge bc$.}
\end{schema}

\begin{strategia}
Ancora una volta riduciamo tutto a numeri non negativi: $a\le b$ significa $b-a\ge0$.
Nel caso (i) anche $c\ge0$, quindi per \ref{ass:ordine-prodotto} il prodotto
$(b-a)\cdot c$ è $\ge0$; sviluppandolo con la proprietà distributiva si trova
$bc-ac\ge0$, che è la tesi. Nel caso (ii) è $-c\ge0$ (\cref{prop:segno-opposto}), e si
ragiona allo stesso modo con $-c$ al posto di $c$: è qui che la disuguaglianza
\emph{cambia verso}.
\end{strategia}

\begin{dimostrazione}
(i) Se $a\le b$ allora per la \cref{prop:ordine-differenza} è $0\le b-a$, da cui, per
\ref{ass:ordine-prodotto} e per la proprietà distributiva (insieme alla regola dei segni,
\cref{prop:segni}),
\[
  0\le (b-a)\cdot c = b\cdot c - a\cdot c,
\]
cioè, ancora per la \cref{prop:ordine-differenza}, $a\cdot c\le b\cdot c$.

(ii) Per le ipotesi e per la \cref{prop:segno-opposto} si ha $0\le b-a$ e $0\le -c$.
Da \ref{ass:ordine-prodotto} si ottiene
\[
  0\le (b-a)\cdot(-c) = -b\cdot c + a\cdot c,
\]
da cui, per la \cref{prop:ordine-differenza}, $b\cdot c\le a\cdot c$.
\end{dimostrazione}

\begin{esempio}[Risolvere una disequazione \emph{citando gli assiomi}]
Risolviamo la disequazione $3-2x\ge 7$ giustificando ogni passaggio.
\begin{align*}
  3-2x\ge7 &\ \sse\ (3-2x)+(-3)\ge 7+(-3) \\
           &\ \sse\ -2x\ge4 \\
           &\ \sse\ (-2x)\cdot\bigl(-\tfrac12\bigr)\le 4\cdot\bigl(-\tfrac12\bigr) \\
           &\ \sse\ x\le-2 .
\end{align*}
Il primo passaggio è l'assioma \ref{ass:ordine-somma} con $c=-3$ (sommare lo stesso
numero ai due membri conserva la disuguaglianza); il terzo è la
\cref{prop:ordine-prodotto}(ii) con $c=-\tfrac12\le0$: il «cambio di verso», che a
scuola si impara come regola mnemonica, è esattamente quella proposizione.
Naturalmente, nel seguito del libro non citeremo più gli assiomi a ogni passaggio: ma è bene sapere che si potrebbe.
\end{esempio}

\begin{esercizio}\label{es:zero-minore-uno}
Dimostrare, a partire dagli assiomi, che $0<1$.
\begin{soluzione}
Per \ref{ass:dicotomia} vale $0\le1$ oppure $1\le0$. Se fosse $1\le0$, per la
\cref{prop:segno-opposto} sarebbe $-1\ge0$, e per \ref{ass:ordine-prodotto}
$(-1)\cdot(-1)\ge0$; ma per la regola dei segni $(-1)\cdot(-1)=1\cdot1=1$, quindi
$1\ge0$. Insieme a $1\le0$, per la proprietà antisimmetrica \ref{ass:antisimmetrica} si
avrebbe $1=0$, contro l'assioma \ref{ass:neutri}. Dunque $0\le1$, e poiché $0\neq1$,
$0<1$. \emph{Scelta logica:} la dicotomia lascia due sole possibilità, e quella sbagliata
si esclude per assurdo; l'idea chiave è che $1$ è un quadrato, $1=(-1)\cdot(-1)$, e un
prodotto di due fattori dello stesso segno non è mai negativo.
\end{soluzione}
\end{esercizio}

\begin{esercizio}\label{es:inverso-positivo}
Dimostrare che se $a>0$ allora $a^{-1}>0$. (\emph{Suggerimento:} se fosse $a^{-1}\le0$,
che cosa si potrebbe dire del prodotto $a\cdot a^{-1}$?)
\begin{soluzione}
Sia $a>0$ e supponiamo per assurdo che sia $a^{-1}\le0$. Allora $-a^{-1}\ge0$
(\cref{prop:segno-opposto}) e, poiché anche $a\ge0$, per \ref{ass:ordine-prodotto} e per
la regola dei segni (\cref{prop:segni})
\[
  0\le a\cdot(-a^{-1})=-(a\cdot a^{-1})=-1 ,
\]
cioè $1\le0$ (ancora per la \cref{prop:segno-opposto}, essendo $-(-1)=1$), contro il
risultato $0<1$ dell'\cref{es:zero-minore-uno}. Dunque non è $a^{-1}\le0$: per la
dicotomia \ref{ass:dicotomia} è $a^{-1}\ge0$ e $a^{-1}\neq0$, cioè $a^{-1}>0$.
\emph{Scelta logica:} la tesi è una disuguaglianza stretta, e la sua negazione
$a^{-1}\le0$ si può moltiplicare per $a\ge0$ grazie a \ref{ass:ordine-prodotto}; il
prodotto $a\cdot a^{-1}=1$ è l'unica informazione che lega $a^{-1}$ ad $a$, e porta a
contraddire $0<1$.
\end{soluzione}
\end{esercizio}
